% NB11b -- Arquitectura del modelo integrado del Wankel: geometria epitrocoidal 2D -> volumen de 3 camaras ->
% termoquimica 0D de dos zonas (EDO en el tiempo) -> fuerzas en los flancos -> par en el eje excentrico.
\documentclass[tikz,border=4mm]{standalone}
\usepackage{amsmath}
\usetikzlibrary{arrows.meta,calc,positioning}
\begin{document}
% Estilo comun de los diagramas TikZ del modulo VCR (NB10b) -- paleta del proyecto
\definecolor{manivela}{HTML}{2A78D6}
\definecolor{inferior}{HTML}{4A3AA7}
\definecolor{control}{HTML}{1BAF7A}
\definecolor{superior}{HTML}{EB6834}
\definecolor{tinta}{HTML}{3A3A38}
\definecolor{metal}{HTML}{C3C2B7}
\definecolor{perdida}{HTML}{E34948}
\definecolor{fuerza}{HTML}{C0392B}
\tikzset{
  barra/.style={line width=3.2pt, line cap=round},
  pin/.style={circle, fill=white, draw=tinta, line width=0.9pt, inner sep=1.6pt},
  fijo/.style={circle, fill=tinta, inner sep=2.6pt},
  vec/.style={-{Stealth[length=2.6mm]}, line width=0.9pt},
  F/.style={-{Stealth[length=3mm]}, line width=1.2pt, fuerza},
  cota/.style={{Stealth[length=1.8mm]}-{Stealth[length=1.8mm]}, line width=0.5pt, tinta},
  lbl/.style={font=\small, text=tinta},
}
\newcommand{\tierra}[1]{\fill[pattern=north east lines, pattern color=tinta] ($(#1)+(-5,-3)$) rectangle ($(#1)+(5,-6)$);
  \draw[tinta, line width=0.8pt] ($(#1)+(-5,-3)$) -- ($(#1)+(5,-3)$);}

\definecolor{mec}{HTML}{2A78D6}\definecolor{ter}{HTML}{E34948}\definecolor{geo}{HTML}{1BAF7A}
\begin{tikzpicture}[x=1cm,y=1cm,
  bloque/.style={draw, line width=1.2pt, rounded corners=4pt, align=left, font=\small, inner sep=6pt, text width=6.4cm},
  info/.style={-{Stealth[length=3mm]}, line width=1.1pt, tinta},
  et/.style={font=\footnotesize, text=tinta, fill=white, inner sep=1.5pt, align=center}]
  \node[bloque, draw=geo, fill=geo!8] (G) at (0,4.6) {%
    \textbf{Geometr\'ia 2D} \hfill {\footnotesize exacta}\\
    estator: epitrocoide $x=R\cos\theta+e\cos3\theta$\\
    rotor: envolvente interior; engranaje $2e{:}3e$\\
    $V_i(\theta)=V_{min}+\tfrac{V_H}{2}\big[1-\sin(\tfrac23\theta+\varphi_i)\big]$, $V_H=3\sqrt3\,eRb$};
  \node[bloque, draw=ter, fill=ter!8] (T) at (4.8,0) {%
    \textbf{Termoqu\'imica 0D} \hfill {\footnotesize EDO en $t=\theta/\omega_e$}\\
    dos zonas $u/b$ por c\'amara (Cantera)\\
    combusti\'on: $\dot m_b$ (Wiebe) desde la chispa\\
    gas sin quemar: iso-octano 143 especies\\
    Woschni con \'area de pared del Wankel};
  \node[bloque, draw=mec, fill=mec!8] (M) at (-4.8,0) {%
    \textbf{Mec\'anica vectorial} \hfill {\footnotesize cuasi-est\'atica}\\
    $\vec F_i=(p_i-p_c)\sqrt3Rb\,\hat n_i$ (por $C$)\\
    $\tau_{eje}=\sum_i\vec e\times\vec F_i$\\
    verificaci\'on: $\tau\,\omega_e=\sum p_i\dot V_i$};
  \node[bloque, draw=tinta, fill=metal!25] (E) at (0,-4.4) {%
    \textbf{Motor} \hfill {\footnotesize 3 c\'amaras desfasadas $360^\circ$, $n$ rotores}\\
    par instant\'aneo y medio, IMEP, BMEP (FMEP medida),\\
    potencia vs. rpm, $C_s$ con volante};
  \draw[info] (G.south east) to[bend left=12] node[et] {$V_i(t),\ \dot V_i(t)$} (T.north);
  \draw[info] (G.south west) to[bend right=12] node[et] {$\vec e(\theta),\ \hat n_i(\theta)$} (M.north);
  \draw[info] (T.west) -- node[et] {$p_i(t)$} (M.east);
  \draw[info] (M.south) to[bend right=12] node[et] {$\tau_{eje}(\theta)$} (E.north west);
  \draw[info] (T.south) to[bend left=12] node[et] {$W_i,\ \dot Q_w$} (E.north east);
\end{tikzpicture}
\end{document}
